\chapter{YOLO, from one grid to YOLO26}

\section{The original idea}

YOLO~\cite{redmon2016yolo} turned detection into a single regression. The image is divided into an $S\times S$ grid.
Each cell predicts a few boxes, each with a confidence, and one set of class probabilities. A convolutional network
reads the whole image once and writes all of these numbers into one output tensor (Figure~\ref{fig:yolov1}).
Because there is no separate region-proposal stage, the detector runs in real time, which is why every later
version kept the name.

\fig{fig_ref_yolov1}{YOLOv1, redrawn after Redmon et al.~\cite{redmon2016yolo}. Top: the network, 24 convolutional
layers followed by two fully connected ones, ending in a $7\times7\times30$ tensor. Bottom: how that tensor is read.
Each grid cell proposes boxes with confidences and a class distribution; their product, after non-maximum
suppression, gives the final detections.}{fig:yolov1}

The weaknesses of this first design shaped everything that followed. One coarse grid means a cell can hold only
one object, so small objects that share a cell are lost. The fully connected head fixes the input size. And
duplicate boxes still have to be removed after the network by non-maximum suppression (NMS), a step that is
neither learned nor parallel.

\section{What changed between versions}

Figure~\ref{fig:lineage} summarises the steps that matter for this project. YOLOv3~\cite{redmon2018yolov3}
predicted at three scales with a feature pyramid~\cite{lin2017fpn}, which is the first answer to small objects.
YOLOv4~\cite{bochkovskiy2020yolov4} and YOLOv5~\cite{jocher2020yolov5} added cross-stage partial backbones,
spatial pyramid pooling and the path-aggregation neck~\cite{liu2018panet}. YOLOv8~\cite{jocher2023yolov8} dropped
anchor boxes, split the head into box and class branches, assigned targets by task alignment~\cite{feng2021tood},
and regressed boxes as distributions (DFL)~\cite{li2020gfl}. YOLOv10~\cite{wang2024yolov10} removed NMS by training
two heads, one that assigns many predictions to each object and one that assigns exactly one, and keeping only the
second at test time.

\widefig{fig_ref_yolo_lineage}{Six generations of YOLO, each reduced to backbone, neck and head, with the change
that defines it. The design converges on a multi-scale pyramid, a path-aggregation neck and an end-to-end head
without NMS.}{fig:lineage}

\section{YOLO26-n, the detector this project builds on}

YOLO26~\cite{yolo26_2026} keeps YOLOv10's end-to-end head and removes the distribution focal loss, regressing box
edges directly. Figure~\ref{fig:yolo26} draws its smallest scale, YOLO26-n, exactly as configured in Ultralytics
8.4.171, with the tensor shapes measured on a 640-pixel input. It has 2.57 million parameters and costs 6.1
GFLOPs.

The backbone halves the resolution five times. Its outputs at strides 8, 16 and 32, called P3, P4 and P5, feed the
neck. The neck first passes information downwards, from the coarse, semantic P5 to the fine P3, by upsampling and
concatenation, then upwards again by strided convolution, so that every output level has both fine detail and
context. Three detection heads read the levels at strides 8, 16 and 32. In training, each head is supervised twice:
a one-to-many branch assigns up to ten predictions to each object, which gives dense gradients, and a one-to-one
branch assigns exactly one, which teaches the network not to produce duplicates. At test time only the one-to-one
branch runs, and the top 300 boxes are the answer, with no NMS.

\widefig{fig_ref_yolo26}{YOLO26-n as used in this project. Blue: backbone; green: the two-way PAN neck; gold: the
end-to-end heads. Numbers left of each block are its layer index; shapes are channels $\times$ height$^2$. The
one-to-many branch is used only in training.}{fig:yolo26}

The four repeated blocks are drawn in Figure~\ref{fig:yoloblocks}. C3k2 is a cross-stage partial block: it splits
the channels, sends one half through a stack of bottlenecks and concatenates the halves again, so gradients have a
short path. SPPF pools the same map three times with a $5\times5$ window to mix context from growing regions.
C2PSA adds self-attention~\cite{vaswani2017attention} at the coarsest level, where the map is small enough for
attention to be cheap.

\fig{fig_ref_yolo26_blocks}{The building blocks of YOLO26: the standard convolution, the cross-stage partial C3k2,
fast spatial pyramid pooling (SPPF), and C2PSA with multi-head attention and a feed-forward layer.}{fig:yoloblocks}
